\begin{frame}{왜 중요도 샘플링이 필요 없는가 --- 5.3절 연결}
  \begin{itemize}
    \item Week 5.3: off-policy MC는 행동/목표 정책의 \textbf{리턴
      분포}가 다르므로 중요도 가중치
      $\rho = \prod_t \pi(a_t\mid s_t)/b(a_t\mid s_t)$로 보정해야.
    \item 그런데 Q-learning은 off-policy임에도
      \textbf{중요도 샘플링을 쓰지 않는다}.
  \end{itemize}
  \vskip0.6em
  \begin{block}{이유}
    Q-learning은 ``목표 정책의 \textbf{리턴 분포}를 샘플로 추정''하는
    것이 아니라, \textbf{벨만 최적방정식의 고정점}을 샘플로 근사하기
    때문. $\max_{a'}Q(s',a')$라는 목표값 자체가 ``목표(탐욕적)
    정책이었다면 다음에 이렇게 갔을 것''을 \textbf{직접} 부트스트랩
   하므로, 행동 정책이 실제로 어떻게 움직였는지에 대한 보정(가중치)이
    구조적으로 불필요.
  \end{block}
  \vskip0.4em
  MC는 ``리턴을 평균''해야 하므로 보정 필요, TD는 ``한 스텝 목표값''
  을 쓰므로 보정 불필요 --- 같은 off-policy 문제를
  \textbf{분포 보정(중요도 샘플링)}과 \textbf{고정점 근사
  (부트스트래핑)}라는 전혀 다른 두 도구로. 이것이 Week 9(DQN)에서
  경험재현 버퍼(다른 정책이 모은 데이터를 재사용)를 자유롭게 쓸 수
  있는 이유이기도 하다.
\end{frame}
